\documentclass{article}
\title{Exocortex}
\author{Ryan Marr}
\date{2025}

\begin{document}
\maketitle

\begin{abstract}
Exocortex began as a Neovim plugin and gradually expanded into an LLM-assisted
knowledge management system. Beyond an expressive, Vim-native interface for
reviewing and accepting LLM-generated code changes, Exocortex organizes
conversations with LLMs as a knowledge graph. That graph is the foundation for
GraphRAG, letting future interactions retrieve relevant context and build on
previous ideas.
\end{abstract}

\section{Motivation}

Conversations with a language model are stored as transcripts: linear, append-only,
and effectively write-only. The reasoning is in there somewhere, but nothing about
the format makes it retrievable, and so the same ground gets re-derived from scratch
in the next session.

The unit worth keeping is not the conversation but the \emph{idea}, along with the
edges connecting it to the ideas it came from.

\section{Components}

\begin{itemize}
  \item \resumeItem{Vim-native review}{An interface for reviewing and accepting
  LLM-generated code changes without leaving the editor, using motions and text
  objects rather than a chat sidebar.}

  \item \resumeItem{Knowledge graph}{Conversations are decomposed into nodes and
  the relationships between them, rather than stored as flat transcripts.}

  \item \resumeItem{GraphRAG retrieval}{The graph backs retrieval for later
  interactions, so a new question is answered with the context of what was already
  worked out.}
\end{itemize}

\section{Status}

Active. Source at
\href{https://github.com/rmarrcode/exocortex.nvim}{github.com/rmarrcode/exocortex.nvim}.

\end{document}
